\section{训练流水线总览：从 SFT 到 Reasoning RL 再到 Agent RL}

\subsection{三阶段训练架构}
课程中将主干流程抽象为三层：监督微调（SFT）奠定基础行为，推理强化学习（Reasoning RL）增强中间思维质量，Agent 强化学习（Agent RL）把能力迁移到环境交互。

\begin{table}[H]
\centering
\begin{tabular}{p{2.7cm}p{4.5cm}p{4.7cm}}
\toprule
\textbf{阶段} & \textbf{目标} & \textbf{典型风险} \\
\midrule
SFT & 学习指令遵循、格式稳定、基础任务完成率 & 过拟合模板，推理深度不足 \\
Reasoning RL & 提升多步推导、策略反思、错误恢复能力 & 奖励设计失真，出现 reward hacking \\
Agent RL & 在真实或模拟环境中优化动作序列与终局成功率 & rollout 成本高，评估噪声大，训练不稳定 \\
\bottomrule
\end{tabular}
\caption{课程中的三阶段训练框架}
\end{table}

\begin{importantbox}{为什么 Reasoning RL 要单独成段}
讲者明确指出，推理强化学习不是 ``顺便做一下'' 的附属步骤，而是连接文本能力与 Agent 行为的关键桥梁。没有稳定的中间推理，Agent 在工具链中的错误会沿时间轴放大，最终在长任务中崩溃。
\end{importantbox}

\subsection{Reasoning RL 的定位与输入输出}
在课程语境下，Reasoning RL 的训练对象不是单一答案，而是 ``思考过程 + 决策路径''。这要求训练样本包含可对齐的中间轨迹，奖励函数也要能区分 ``虽然答对但过程不可复用'' 与 ``过程稳定可迁移'' 的区别。

\begin{warningbox}{常见误区：把推理 token 量当作推理质量}
更长的 chain-of-thought 不等于更好的策略。课程强调真正的收益来自 \textbf{有结构的推理行为}，例如何时回退、何时重计划、何时调用工具，而不是单纯延长输出长度。
\end{warningbox}

\subsection{本章小结}
SFT、Reasoning RL、Agent RL 构成了一个由浅入深的能力迁移链。Reasoning RL 的存在不是可选项，而是把 ``会说'' 变成 ``会做'' 的中介层。


